    \hspace{3mm}

\centerline{\bf \Huge Теория по многочленам}

\begin{flushright}
    \scriptsize{- Ха-ха, ха-ха-ха! Шо ты, стеинг элайв?\\Жотаро Токийский}
\end{flushright}


\textbf{\Large Глава 1. Что такое многочлен?}

Многочленом называют формальное выражение вида $a_nx^n + a_{n-1}x^{n-1}+\dotsc+a_1x+a_0$, где $a_0, a_1, a_2, \dotsc , a_n$ называются коэффициентами этого многочлена, а $x$ это переменная, от которой наш многочлен зависит.

Пример многочленов:

$x^3+x+1$

$x$

$1 \ - $ это тоже многочлен, просто нулевой степени

$ax^2 + bx + c \ -$ квадратный трехчлен, где $a, b, c$ неизвестные коэффициенты

Мы будем рассматривать только многочлены от одной переменной, большего нам не надо. Многочлены обозначаются латинскими заглавными буквами(самые классические это $P(x)$ и $Q(x)$). В выражение $P(x)$ вместо переменной $x$ можно подставлять конкретные числа, как в обычную функцию(многочлен это по сути и есть функция от переменной $x$). Например, если $P(x) = x^2+x+1$, то $P(2) = 2^2 + 2 + 1 = 7$.

\begin{itemize}
    \item Многочлен первой степени $P(x) = ax+b$ называется линейным двучленом или биномом.
    \item Многочлен второй степени $P(x) = ax^2 + bx+c$ называется квадратным трёхчленом
    \item Многочлен третьей степени называется кубическим
\end{itemize}

\textbf{\Large Глава 2. Оказывается многочлены можно делить друг на друга}

Да, можно поделить один многочлен на другой с остатком и делается точно так же, как и с обычными числами $-$ уголком! Я человек очень ленивый и набирать все эти деления уголком в LaTeX мне лень. Поэтому держите \href{https://www.youtube.com/watch?v=yhJ9xeo_W8c}{ссылку}, где добрая тетенька очень доступно поясняет, как же делить многочлены друг на друга. Единственное добавляю от себя, что \textbf{степень остатка всегда строго меньше чем степень делителя!!!}

\newpage

\hspace{3mm}

\textbf{\Large Глава 3. Теорема Безу}

Одно из самых важных утверждений про многочлены это теорема Безу. Давайте сформулируем и сразу же докажем её.

\textit{Теорема Безу.} Многочлен $P(x)$ дает остаток $P(a)$ при делении на $x-a$.

Доказательство:

\noindentПоделим $P(x)$ на $x-a$ с остатком. Пусть получилось частное $Q(x)$ и остаток $R(x)$. Тогда $P(x)$ представляется как $P(x) = Q(x)(x-a)+R(x)$, но степень $R(x)$ строго меньше чем степень делителя, то есть чем степень многочлена $x-a$, которая равна 1. Значит степень многочлена $R(x)$ нулевая и это просто константа, которую мы обозначим за $r$. Тогда получаем, что $P(x) = Q(x)(x-a)+r$. Подставим $x=a$ и получим, что $P(a)=r$ то есть остаток. Что и требовалось доказать!

\textbf{Примеры применения теоремы Безу}

\problem{1} Найти остаток от деления многочлена $x^{1001} + x^{13} + x^{11} + x^{7}$ при делении на $x+1$.

\textit{Решение: }Конечно, мы бы не ослили это все вручную делить в столбик и наверняка где-нибудь по дороге ошиблись бы. Поэтому изящно применяем теорему Безу и не забываем следить за знаками, так как делим мы на $x+1$, а не на $x-1$. Поэтому подставлять нужно вместо $x$ именно -1. Получаем остаток $(-1)^{1001}+(-1)^{13}+(-1)^{11}+(-1)^{7}=-4$

\problem{2} Докажите, что многочлен $P(x) = (x+1)^6 - x^6 - 2x-1$ делится на $x(x+1)(2x+1)$

\textit{Решение: }Давайте сперва покажем, что наш многочлен делится на $x$ и $x+1$. Для этого достаточно проверить, что $P(0) = P(-1) = 0$, что правда. Но как быть с делимость на $2x+1$? Теорема Безу нам говорит только про делимость на $x-a$ а тут откуда-то коэффициент при $x$. Что же с ним делать? Давайте просто вынесем его: $2x+1=2(x+\dfrac{1}{2})$. Делимость на константу для многочленов очевидна, любой многочлен можно безболезнено поделить на ненулевую константу(если нет ограничения на коэффициенты). Теперь уже спокойно применяем теорему Безу, подставляем $-\dfrac{1}{2}$ и кайфуем.


\newpage

\hspace{3mm}

\textbf{\Large Глава 4. Целочисленный Безу}

Прежде чем говорить о теореме Безу для многочленов с целыми коэффициентами, вспоминим очень полезные формулы сокращенного умножения, которые \textbf{обязательно} знать каждому олимпиаднику.

\problem{1} Формула разности степеней.

$a^k - b^k = (a-b)(a^{k-1}+a^{k-2}b+a^{k-3}b^2+\dotsc+a^2b^{k-3}+ab^{k-2}+b^{k-1})$


\problem{2} Формула суммы степеней. Её можно использовать только для нечётных степеней.

$a^{2k+1} + b^{2k+1} = (a+b)(a^{2k-1}-a^{2k-2}b+a^{2k-3}b^2-\dotsc+a^2b^{k-3}-ab^{k-2}+b^{k-1})$


\noindentЗапомнить формулы довольно легко. В маленькой скобке такой же знак как и был в начале, а в большой скобке все плюсы, если был минус и знаки чередуются, если был плюс. Из этих двух формул можно выделить сразу два следствия:

Следствие 1. Число $a^k - b^k$ делится на $a-b$

Следствие 2. Число $a^{2k+1}+b^{2k+1}$ делится на $a+b$

\hspace{3mm}

Теперь можем сформулировать и доказать теорему Безу для целочисленных многочленов.

\textit{Теорема Безу для многочленов с целыми коэффициентами.} Пусть $P(x)$ многочлен с целыми коэффициентами и числа $a, b$ тоже целые. Тогда верно, что $P(a)-P(b)$ делится на $a-b$.

Доказательство:

\noindentПусть $P(x) = a_nx^n + a_{n-1}x^{n-1}+\dotsc+a_1x+a_0$, где $a_n, a_{n-1}, \dotsc , a_1, a_0$ целые числа. Посмотрим, что такое $P(a)-P(b)$

$P(a)-P(b) = a_na^n + a_{n-1}a^{n-1}+\dotsc+a_1a+a_0 - (a_nb^n + a_{n-1}b^{n-1}+\dotsc+a_1b+a_0) = a_n(a^n-b^n) + a_{n-1}(a^{n-1}-b^{n-1})+\dotsc+a_1(a-b)+(a_0-a_0)$

Вспоминаем наше следствие из формулы сокращенного умножения и замечаем, что каждое слагаемое делится на $a-b$, а значит и все выражение. Ч.Т.Д.

